\documentclass[12pt,a4paper]{article}
\usepackage[margin=2cm]{geometry}
\usepackage{booktabs}
\usepackage{xcolor}
\usepackage{amsmath}
\usepackage{enumitem}
\usepackage{hyperref}
\hypersetup{colorlinks=true, linkcolor=blue, urlcolor=blue}

\definecolor{goodgreen}{RGB}{0,120,0}
\definecolor{badred}{RGB}{180,0,0}

\title{Binary vs Multi-Class Classification:\\Why 6-Class Accuracy Drops and How to Fix It}
\author{Landslide Identification Using ML}
\date{March 2026}

\begin{document}
\maketitle

% ============================================================
\section{The Two Experiments}
% ============================================================

We ran the same EfficientNet-B3 model on the same HR-GLDD dataset with two configurations:

\begin{itemize}
    \item \textbf{Binary (NUM\_CLASSES = 2):} Landslide vs Non-Landslide
    \item \textbf{Multi-class (NUM\_CLASSES = 6):} Non-Landslide + 5 landslide subtypes
\end{itemize}

% ============================================================
\section{Results Comparison}
% ============================================================

\begin{table}[h]
    \centering
    \renewcommand{\arraystretch}{1.3}
    \begin{tabular}{lcc}
        \toprule
        \textbf{Metric} & \textbf{Binary (2-class)} & \textbf{Multi-class (6-class)} \\
        \midrule
        Accuracy        & \textcolor{goodgreen}{\textbf{84.55\%}} & \textcolor{badred}{41.34\%} \\
        Precision (macro) & \textcolor{goodgreen}{\textbf{71.50\%}} & \textcolor{badred}{22.42\%} \\
        Recall (macro)    & \textcolor{goodgreen}{\textbf{81.04\%}} & \textcolor{badred}{22.20\%} \\
        F1-Score (macro)  & \textcolor{goodgreen}{\textbf{75.98\%}} & \textcolor{badred}{20.18\%} \\
        ROC-AUC (macro)   & \textcolor{goodgreen}{\textbf{91.50\%}} & \textcolor{badred}{64.81\%} \\
        \bottomrule
    \end{tabular}
    \caption{Binary vs Multi-class results on the same HR-GLDD test set (699 images).}
\end{table}

\noindent\textbf{The drop:} Accuracy fell from 84.55\% to 41.34\% --- a \textbf{43 percentage point} drop. F1 fell from 75.98\% to 20.18\%.

\subsection*{Per-Class Breakdown (6-class)}

\begin{table}[h]
    \centering
    \renewcommand{\arraystretch}{1.2}
    \small
    \begin{tabular}{lcccc}
        \toprule
        \textbf{Class} & \textbf{Training Samples} & \textbf{Precision} & \textbf{Recall} & \textbf{F1} \\
        \midrule
        Non-landslide       & 1,574 & 92.73\% & 52.25\% & 66.84\% \\
        Rockfall             & $\sim$120 & 6.60\% & 20.59\% & 10.00\% \\
        Mudflow              & $\sim$120 & 8.89\% & 18.18\% & 11.94\% \\
        Debris flow          & $\sim$130 & 11.39\% & 18.37\% & 14.06\% \\
        Rotational slide     & $\sim$120 & 3.41\% & 7.14\% & 4.62\% \\
        Translational slide  & $\sim$126 & 11.48\% & 16.67\% & 13.59\% \\
        \bottomrule
    \end{tabular}
    \caption{6-class per-class metrics. Non-landslide dominates; landslide subtypes are near random.}
\end{table}

\noindent\textbf{Key observation:} Non-landslide class gets 92.73\% precision because it has 1,574 training images. All landslide subtypes are below 12\% precision --- the model essentially cannot distinguish between them.

% ============================================================
\section{Why Does 6-Class Accuracy Drop So Much?}
% ============================================================

There are \textbf{three root causes}, each compounding the others:

\subsection{Cause 1: Insufficient Training Samples Per Class}

\begin{table}[h]
    \centering
    \begin{tabular}{lcc}
        \toprule
        \textbf{Setup} & \textbf{Smallest Class} & \textbf{Enough?} \\
        \midrule
        Binary   & 616 (all landslides combined) & Marginal but workable \\
        6-class  & $\sim$120 (per subtype) & \textcolor{badred}{\textbf{No --- far too few}} \\
        \midrule
        ImageNet benchmark & 1,200+ per class & Standard minimum \\
        \bottomrule
    \end{tabular}
\end{table}

Deep learning models need \textbf{hundreds to thousands} of images per class to learn reliable patterns. With only $\sim$120 images per landslide subtype, the model does not have enough examples to learn what makes a rockfall look different from a mudflow.

\textbf{Simple analogy:} Imagine showing someone 120 photos of cats and 120 photos of dogs, then asking them to also distinguish breed. They could tell cat vs dog (binary) easily, but distinguishing Labrador vs Golden Retriever from just 120 photos each would be much harder.

\subsection{Cause 2: High Visual Similarity Between Subtypes}

From satellite view at 128$\times$128 pixels, landslide subtypes look extremely similar:

\begin{itemize}
    \item \textbf{Mudflow:} Brown soil flowing down a slope
    \item \textbf{Debris flow:} Brown soil and rocks flowing down a slope
    \item \textbf{Rotational slide:} Curved scar where soil moved downhill
    \item \textbf{Translational slide:} Straight scar where soil moved downhill
    \item \textbf{Rockfall:} Rocky debris at base of a cliff
\end{itemize}

The visual difference between a mudflow and a debris flow from overhead is often \textbf{invisible in RGB satellite imagery}. Even expert geologists need:
\begin{itemize}
    \item Ground-level photos (not just satellite)
    \item Geological survey data (soil composition)
    \item Movement speed data (slow vs fast)
    \item Terrain slope and elevation profiles
\end{itemize}

A 128$\times$128 RGB satellite image simply does not contain enough information to make these distinctions reliably.

\subsection{Cause 3: Severe Class Imbalance (13:1 ratio)}

\begin{itemize}
    \item Non-landslide: \textbf{1,574} training images (72\%)
    \item Each landslide subtype: \textbf{$\sim$120} training images ($\sim$5.5\% each)
    \item Imbalance ratio: \textbf{13:1} (non-landslide vs any subtype)
\end{itemize}

The model learns that predicting ``non-landslide'' for everything gives $\sim$70\% accuracy with zero effort. There is not enough signal in the rare classes to overcome this bias, even with Focal Loss and WeightedRandomSampler.

% ============================================================
\section{What Would Be Needed to Achieve High 6-Class Accuracy?}
% ============================================================

\subsection{Approach 1: More Data (Most Important)}

\begin{table}[h]
    \centering
    \renewcommand{\arraystretch}{1.2}
    \begin{tabular}{lccl}
        \toprule
        \textbf{Dataset} & \textbf{Images} & \textbf{Subtype Labels?} & \textbf{Problem} \\
        \midrule
        HR-GLDD (ours) & 3,439 & Yes (5 subtypes) & Too few per class \\
        CAS Landslide & 20,865 & No --- binary only & Cannot use for subtypes \\
        Landslide4Sense & 3,799 & No --- binary only & Cannot use for subtypes \\
        IEEE DataPort & 770 & Partial (3 types) & Missing 2 subtypes \\
        DMLD & 990 & No --- binary only & Cannot use for subtypes \\
        \bottomrule
    \end{tabular}
    \caption{Available public datasets. Almost none have subtype labels.}
\end{table}

\textbf{The core problem:} Almost no public dataset labels landslide \textit{subtypes}. Most only label ``landslide vs non-landslide.'' To build a proper 6-class dataset, you would need:

\begin{itemize}
    \item \textbf{500--1,000+ images per subtype} (minimum)
    \item Images labeled by a \textbf{domain expert geologist}
    \item Higher resolution imagery (\textbf{$<$1m/pixel}) to see subtype-specific features
    \item Multiple data sources: RGB + \textbf{NIR band} (for NDVI) + \textbf{DEM} (elevation/slope)
\end{itemize}

\textbf{Estimated effort:} 2--4 months of data collection and manual annotation.

\subsection{Approach 2: Hierarchical Classification (Best Practical Solution)}

Instead of one model doing all 6 classes, use \textbf{two stages}:

\begin{enumerate}
    \item \textbf{Stage 1 --- Binary Detection:} Landslide or not?
    \begin{itemize}
        \item Uses the current ensemble (EfficientNet-B3 + ViT-B/16)
        \item Already achieves \textbf{84.55\% accuracy}
        \item Filters out 488 non-landslide images, leaving only 211 positives
    \end{itemize}

    \item \textbf{Stage 2 --- Subtype Classification:} What type of landslide?
    \begin{itemize}
        \item A separate model trained ONLY on the 211 landslide images
        \item No non-landslide class to dominate --- all 5 classes are balanced
        \item Combined with HMM type estimates for cross-validation
    \end{itemize}
\end{enumerate}

\textbf{Why this works better:}
\begin{itemize}
    \item Stage 1 stays at 84.55\% (proven)
    \item Stage 2 only needs to distinguish 5 subtypes from each other (easier than 6 classes with a dominant majority class)
    \item The 13:1 imbalance problem disappears because non-landslide is removed
\end{itemize}

\subsection{Approach 3: Better Input Features}

\begin{itemize}
    \item \textbf{NIR band:} HR-GLDD actually has a 4th NIR channel. Using it for NDVI (vegetation index) would help distinguish subtypes based on vegetation damage patterns.
    \item \textbf{DEM + Slope:} Elevation and terrain slope data helps because rockfalls occur on steep cliffs while rotational slides occur on moderate slopes.
    \item \textbf{Higher resolution:} 128$\times$128 at 3m/pixel only covers a 384m$\times$384m area. Fine details that distinguish subtypes are lost.
\end{itemize}

\subsection{Approach 4: Advanced Augmentation}

\begin{itemize}
    \item \textbf{Mixup / CutMix:} Synthesize training images by blending existing ones
    \item \textbf{Generative augmentation:} Use diffusion models to generate synthetic landslide subtype images
    \item \textbf{Cross-dataset transfer:} Pretrain on CAS Landslide (20K images, binary) then fine-tune on HR-GLDD (subtypes)
\end{itemize}

% ============================================================
\section{Summary: What We Recommend}
% ============================================================

\begin{table}[h]
    \centering
    \renewcommand{\arraystretch}{1.3}
    \begin{tabular}{lccl}
        \toprule
        \textbf{Approach} & \textbf{Effort} & \textbf{Expected Impact} & \textbf{Timeline} \\
        \midrule
        More labeled data & High & +20--30\% & 2--4 months \\
        Hierarchical classification & Medium & Binary stays 84\% & 2--3 weeks \\
        NIR + DEM features & Medium & +5--10\% & 1--2 weeks \\
        Advanced augmentation & Low & +3--5\% & 1 week \\
        \bottomrule
    \end{tabular}
\end{table}

\vspace{0.3cm}
\noindent\fcolorbox{black}{green!10}{%
\parbox{\dimexpr\linewidth-2\fboxsep-2\fboxrule}{%
\textbf{Bottom line:} The binary model (84.55\%) is the project's core strength. The 6-class experiment is an honest scientific investigation that reveals a real limitation and proposes concrete solutions. The recommended path forward is hierarchical classification, which keeps the binary detection at 84.55\% while adding subtype classification as a second stage.
}}

% ============================================================
\section{How to Explain This to Your Guide}
% ============================================================

If your guide asks ``Why is 6-class accuracy so low?'', say:

\begin{quote}
``The 6-class experiment dropped from 84.55\% to about 41\% for three reasons: first, we only have about 120 training images per landslide subtype, which is far below the minimum needed for deep learning. Second, from satellite view, mudflows and debris flows look almost identical --- even geologists need ground-level data to distinguish them. Third, the 13:1 class imbalance means the model gets rewarded for always predicting non-landslide.

The solution is hierarchical classification: first detect landslide vs non-landslide at 84.55\% accuracy, then classify the subtype only on detected landslides. This removes the imbalance problem and focuses the subtype model on a balanced, easier task. We documented this as future work.''
\end{quote}

This answer shows you understand the \textbf{problem}, the \textbf{cause}, and the \textbf{solution} --- which is exactly what a guide wants to hear.

\end{document}
